% Source: https://artofproblemsolving.com/community/c3h2193077_seesaw_puzzle
% Stable identity: aops-topic-2193077-q01
There are $24$ possibilities (one of $12$ men, heavier or lighter), and three uses of the seesaw have $3^3=27$ possible outcomes, so a strategy is not ruled out; here is one. Number the men $1$ to $12$, and call the man with the different weight the odd man.

\textbf{First weighing: $1,2,3,4$ against $5,6,7,8$.}

\emph{(i) Balance.} Men $1$-$8$ are normal and the odd man is among $9,10,11,12$. Second weighing: $9,10,11$ against $1,2,3$.
\begin{itemize}
\item Balance: the odd man is $12$. Third weighing: $12$ against $1$; the seesaw shows whether $12$ is heavier or lighter (it cannot balance).
\item $9,10,11$ heavier: the odd man is one of them and is heavier. Third weighing: $9$ against $10$; the heavier one is the odd man, and if they balance it is $11$.
\item $9,10,11$ lighter: the same with ``lighter''.
\end{itemize}

\emph{(ii) $1,2,3,4$ heavier.} Then either one of $1,2,3,4$ is heavy or one of $5,6,7,8$ is light, and $9$-$12$ are normal. Second weighing: $1,2,5$ against $3,6,9$.
\begin{itemize}
\item Balance: the odd man is $4$ (heavy), $7$ (light) or $8$ (light). Third weighing: $7$ against $8$; the lighter one is the odd man and is light, and if they balance, $4$ is the odd man and is heavy.
\item $1,2,5$ heavier: either $1$ or $2$ is heavy, or $6$ is light ($5$ cannot be light on the heavier side, and $3$ cannot be heavy on the lighter side). Third weighing: $1$ against $2$; the heavier one is the odd man and is heavy, and if they balance, $6$ is the odd man and is light.
\item $3,6,9$ heavier: either $3$ is heavy or $5$ is light. Third weighing: $3$ against $9$ (a normal man); if $3$ is heavier, $3$ is the odd man and is heavy, and if they balance, $5$ is the odd man and is light.
\end{itemize}

\emph{(iii) $5,6,7,8$ heavier.} This is case (ii) with the roles of ``heavier'' and ``lighter'' exchanged: use exactly the same second and third weighings and swap the words heavy and light in every conclusion.

In every case the odd man and whether he is heavier or lighter are found with three weighings.
